% !TEX TS-program = xelatex
% !TEX encoding = UTF-8 Unicode
% !Mode:: "TeX:UTF-8"

\documentclass{resume}
\usepackage{zh_CN-Adobefonts_external} % Simplified Chinese Support using external fonts (./fonts/zh_CN-Adobe/)
%\usepackage{zh_CN-Adobefonts_internal} % Simplified Chinese Support using system fonts
\usepackage{linespacing_fix} % disable extra space before next section
\usepackage{cite}

\begin{document}
\pagenumbering{gobble} % suppress displaying page number

\name{陈耀}

% {E-mail}{mobilephone}{homepage}
% be careful of _ in emaill address
\contactInfo{cy8530@sjtu.edu.cn} {(+86) 135-9419-5162}{意向岗位：大模型与Agent方向}
% {E-mail}{mobilephone}
% keep the last empty braces!
%\contactInfo{xxx@yuanbin.me}{(+86) 131-221-87xxx}{}

\section{个人总结}
\begin{itemize}
    \item \textbf{聚焦AI前沿技术}：专注于深度学习领域，熟悉CNN、Transformer等主流网络架构；紧跟大模型（LLM）技术发展趋势，具备大模型应用开发及微调实战经验，了解AI Agent及RAG技术应用。
    \item \textbf{自驱力与工程能力}：拥有开阔的技术视野，善于利用开源社区资源快速掌握新框架与新工具；性格开朗乐观，具备强烈的责任心与团队协作意识，渴望在挑战性环境中通过解决实际问题实现快速成长。
\end{itemize}

% \section{\faGraduationCap\ 教育背景}
\section{教育背景}
\datedsubsection{\textbf{上海交通大学},智能感知,\textit{在读本科生}}{2023.9 - 至今}
\ \textbf{排名19/63(前30\%)，国家C等奖学金}，预计2027年6月毕业

% \section{\faCogs\ IT 技能}
\section{技术能力}
% increase linespacing [parsep=0.5ex]
\begin{itemize}[parsep=0.2ex]
  \item \textbf{编程语言}: Python, C++
  \item \textbf{深度学习与数据}: PyTorch, Pandas, NumPy, Hugging Face Transformers
  \item \textbf{大模型应用开发}:
  \begin{itemize}
      \item \textbf{框架与SDK}: 使用 Dashscope、OpenAI SDK 进行模型调用；了解 LlamaIndex, LangChain 及 AgentScope 框架，具备构建 RAG 系统的基础能力。
      \item \textbf{Agent与评估}: 理解 Agent 规划与工具使用机制，掌握上下文工程技巧；了解使用 Ragas 对 RAG 系统进行效果评估。
  \end{itemize}
  \item \textbf{工程与工具}: Linux环境开发, Git, Docker基础, MySQL
\end{itemize}

% \end{itemize}

\section{项目经历}

\datedsubsection{\textbf{基于YOLOv11与对比学习的巨核细胞检测与分类系统} | 瑞金医院合作项目}{2025.09 - 至今}
\begin{itemize}
    \item \textbf{项目背景}：针对传统病理诊断依赖专家人工镜检，存在标注主观性强、耗时耗力且难以规模化等痛点，与瑞金医院合作开展医学图像自动化分析研究。
    \item \textbf{算法创新}：基于 YOLOv11 框架，引入 \textbf{对比学习（Contrastive Learning）} 策略增强特征提取能力，并融合\textbf{先验概率分布}优化分类边界，实现了对三种巨核细胞亚型的高精度检测与分类。
    \item \textbf{成果产出}：项目成果已整理成论文投稿至 \textbf{CCF-C类会议}（\textbf{一作}），目前处于在投状态；代码实现了全流程自动化训练与推理。
\end{itemize}

\datedsubsection{\textbf{基于RAG与AgentScope的智能问答系统全链路实践} | 个人项目（阿里云ACP认证课程）}{2026.04 - 2026.05}
\begin{itemize}
    \item \textbf{系统性实践}：基于阿里云ACP认证课程，完整复现并实践了智能问答系统的标准开发流程，涵盖从文档解析、Embedding向量化、RAG构建、AgentScope集成到Ragas评估的全链路。基于 \textbf{AgentScope} 框架构建 \textbf{ReAct Agent}，并通过 \textbf{MCP (Model Context Protocol)} 协议管理外部工具，实现了 \textbf{Function Calling} 与 \textbf{Skill} 的动态挂载。
    \item \textbf{RAG与上下文工程}：基于 \textbf{LlamaIndex} 搭建 RAG 管道，实施了完整的 \textbf{上下文工程}。通过对比 \textbf{Token/Sentence/Markdown} 等切片策略，以及不同\textbf{embedding}模型，并引入 \textbf{HyDE}（假设文档嵌入）与 \textbf{Rerank} 重排序模型等方法，有效解决了检索召回率低等问题，显著提升了问答准确度。
    \item \textbf{量化评测与反思}：引入 \textbf{Ragas} 框架构建自动化评测体系，利用 \textbf{Answer Correctness} 和 \textbf{Context Recall} 指标驱动迭代；设计 \textbf{Reflection (反思)} 机制，通过独立审查 Agent 进行多轮自我纠错，显著抑制了模型幻觉，提升了系统在专业场景下的可靠性。
\end{itemize}

\datedsubsection{\textbf{基于Spark与Docker的大数据推荐系统原型} | 课程项目/个人项目}{2026.04 - 2026.05}
\begin{itemize}
    \item \textbf{环境搭建}：作为项目负责人，基于 Docker 搭建了包含 \textbf{HDFS、Spark、Hive} 的伪分布式集群环境，构建了端到端的数据处理流水线。
    \item \textbf{数据工程}：使用 Python 脚本生成模拟用户行为数据，利用 Hive 构建数仓进行结构化存储，并通过 \textbf{PySpark} 实现 \textbf{ALS 协同过滤算法} 训练推荐模型。
    \item \textbf{链路打通}：解决了异构数据源类型转换难题，探索了 HBase 集成方案，实现了离线计算结果向在线存储的流转，模型最终持久化至 HDFS。
\end{itemize}


% \section{\faHeartO\ 项目/作品摘要}
% \section{项目/作品摘要}
% \datedline{\textit{An Integrated Version of Security Monitor Vis System}, https://hijiangtao.github.io/ss-vis-component/ }{}
% \datedline{\textit{Dark-Tech}, https://github.com/hijiangtao/dark-tech/ }{}
% \datedline{\textit{融合社交网络数据挖掘的电视节目可视化分析系统}, https://hijiangtao.github.io/variety-show-hot-spot-vis/}{}
% \datedline{\textit{LeetCodeOJ Solutions}, https://github.com/hijiangtao/LeetCodeOJ}{}
% \datedline{\textit{Info-Vis}, https://github.com/ISCAS-VIS/infovis-ucas}{}


%% Reference
%\newpage
%\bibliographystyle{IEEETran}
%\bibliography{mycite}
\end{document}